\begin{ex}
Một hình tam giác có độ dài ba cạnh lần lượt là 12 cm, 15 cm và 18 cm. Tính chu vi hình tam giác đó theo đơn vị xăng-ti-mét.
\shortans{45}
\loigiai{
$12 + 15 + 18 = 45$ cm.
}
\end{ex}

\begin{ex}
Một hình tứ giác có độ dài các cạnh lần lượt là 10 cm, 20 cm, 15 cm và 25 cm. Chu vi hình tứ giác đó là bao nhiêu xăng-ti-mét?
\shortans{70}
\loigiai{
$10 + 20 + 15 + 25 = 70$ cm.
}
\end{ex}

\begin{ex}
Một đoạn thẳng AB dài 15 cm, đoạn thẳng CD dài hơn đoạn AB là 7 cm. Hỏi đoạn thẳng CD dài bao nhiêu xăng-ti-mét?
\shortans{22}
\loigiai{
$15 + 7 = 22$ cm.
}
\end{ex}

\begin{ex}
Một sợi dây dài 80 cm, người ta cắt đi 35 cm. Hỏi đoạn dây còn lại dài bao nhiêu xăng-ti-mét?
\shortans{45}
\loigiai{
$80 - 35 = 45$ cm.
}
\end{ex}

\begin{ex}
Đổi: $1\text{ m} = \dots\text{ cm}$. Điền số thích hợp.
\shortans{100}
\loigiai{
$1\text{ m} = 100\text{ cm}$.
}
\end{ex}

\begin{ex}
Đổi: $1\text{ m} = \dots\text{ dm}$. Điền số thích hợp.
\shortans{10}
\loigiai{
$1\text{ m} = 10\text{ dm}$.
}
\end{ex}

\begin{ex}
Đổi: $3\text{ m } 20\text{ cm} = \dots\text{ cm}$. Điền số thích hợp.
\shortans{320}
\loigiai{
$3\text{ m } 20\text{ cm} = 300 + 20 = 320\text{ cm}$.
}
\end{ex}

\begin{ex}
Đổi: $5\text{ dm } 4\text{ cm} = \dots\text{ cm}$. Điền số thích hợp.
\shortans{54}
\loigiai{
$5\text{ dm } 4\text{ cm} = 50 + 4 = 54\text{ cm}$.
}
\end{ex}

\begin{ex}
Một can đựng 10 lít nước, một xô đựng 5 lít nước. Hỏi cả can và xô đựng bao nhiêu lít nước?
\shortans{15}
\loigiai{
$10 + 5 = 15$ lít.
}
\end{ex}

\begin{ex}
Mẹ mua 25 kg gạo nếp và 30 kg gạo tẻ. Hỏi mẹ mua tất cả bao nhiêu ki-lô-gam gạo?
\shortans{55}
\loigiai{
$25 + 30 = 55$ kg.
}
\end{ex}

\begin{ex}
Một túi đường nặng 2 kg. Hỏi 5 túi đường như thế nặng tất cả bao nhiêu ki-lô-gam?
\shortans{10}
\loigiai{
$2 \times 5 = 10$ kg.
}
\end{ex}

\begin{ex}
Một thùng chứa 40 lít dầu, người ta đã lấy ra 16 lít. Hỏi trong thùng còn lại bao nhiêu lít dầu?
\shortans{24}
\loigiai{
$40 - 16 = 24$ lít.
}
\end{ex}

\begin{ex}
Hình vuông có cạnh dài 5 cm. Tính tổng độ dài 4 cạnh của hình vuông đó theo đơn vị xăng-ti-mét.
\shortans{20}
\loigiai{
$5 \times 4 = 20$ cm.
}
\end{ex}

\begin{ex}
Hình chữ nhật có chiều dài 8 cm, chiều rộng 5 cm. Nửa chu vi hình chữ nhật đó là bao nhiêu xăng-ti-mét?
\shortans{13}
\loigiai{
Nửa chu vi = chiều dài + chiều rộng = $8 + 5 = 13$ cm.
}
\end{ex}

\begin{ex}
Có bao nhiêu cạnh trong một hình tứ giác?
\shortans{4}
\loigiai{
Hình tứ giác có 4 cạnh.
}
\end{ex}

\begin{ex}
Có bao nhiêu đỉnh trong một hình tam giác?
\shortans{3}
\loigiai{
Hình tam giác có 3 đỉnh.
}
\end{ex}

\begin{ex}
Một khối lập phương có bao nhiêu mặt?
\shortans{6}
\loigiai{
Khối lập phương có 6 mặt.
}
\end{ex}

\begin{ex}
Một khối lập phương có bao nhiêu đỉnh?
\shortans{8}
\loigiai{
Khối lập phương có 8 đỉnh.
}
\end{ex}

\begin{ex}
Một khối hộp chữ nhật có bao nhiêu cạnh?
\shortans{12}
\loigiai{
Khối hộp chữ nhật có 12 cạnh.
}
\end{ex}

\begin{ex}
Đường gấp khúc gồm 3 đoạn thẳng có độ dài lần lượt là 14 cm, 16 cm và 20 cm. Tính độ dài đường gấp khúc đó theo đơn vị xăng-ti-mét.
\shortans{50}
\loigiai{
$14 + 16 + 20 = 50$ cm.
}
\end{ex}

\begin{ex}
Đoạn thẳng MN dài 40 cm. Điểm P là trung điểm của MN. Hỏi đoạn thẳng MP dài bao nhiêu xăng-ti-mét?
\shortans{20}
\loigiai{
$40 : 2 = 20$ cm.
}
\end{ex}

\begin{ex}
Rót hết 18 lít mật ong vào các can loại 2 lít. Hỏi rót được bao nhiêu can mật ong?
\shortans{9}
\loigiai{
$18 : 2 = 9$ can.
}
\end{ex}

\begin{ex}
Có 3 gói kẹo, mỗi gói cân nặng 200 g. Hỏi cả 3 gói kẹo cân nặng bao nhiêu gam?
\shortans{600}
\loigiai{
$200 \times 3 = 600$ g.
}
\end{ex}

\begin{ex}
Một chiếc thước kẻ dài 30 cm, bút chì dài 18 cm. Chiếc thước kẻ dài hơn bút chì bao nhiêu xăng-ti-mét?
\shortans{12}
\loigiai{
$30 - 18 = 12$ cm.
}
\end{ex}

\begin{ex}
Chu vi hình tam giác đều có cạnh 7 cm là bao nhiêu xăng-ti-mét?
\shortans{21}
\loigiai{
$7 \times 3 = 21$ cm.
}
\end{ex}
